\documentclass[11pt,a4paper]{article}

\usepackage[margin=1in]{geometry}
\usepackage[T1]{fontenc}
\usepackage[utf8]{inputenc}
\usepackage{lmodern}
\usepackage{hyperref}
\usepackage{xcolor}
\hypersetup{colorlinks=true,linkcolor=blue,urlcolor=blue,citecolor=blue}

\title{Detailed Approach and Evaluation of the LangGraph Agentic Application}
\author{}
\date{}

\begin{document}

\maketitle

\section{Approach}

The LangGraph agentic application was populated through a prompt-driven synthesis process in which IBM BOB interpreted the user request and turned mined process knowledge into an executable application design. The initiating instruction that framed this work was the following: ``The first section should elaborate how the langraph agentic app was populated, initiated by the below prompt, with a focus on the role-specific agents that were populated. Summarize how each agent was specified in the app accordingly. The second section should describe the two issues that the app was tested with (\#1 and \#2) as a kind of a functional testing of the app. Explain the content of the two issue and how they were interpreted by the app, with routing to the relevant agent and LLM based analysis and derivation of a corresponding relevant action. Please write an article style paragraph within minimal use of itemization, and overall not more than 2 pages, 1 for each section.'' In this process, IBM BOB employed a mixture of frontier and specialized language models for interpretation, drawing on Anthropic Claude, Mistral open source models, IBM Granite, and additional specialized fine-tuned models. The app was therefore not handcrafted from scratch in isolation, but generated by combining prompt interpretation with artifacts mined from the repository process, including OCDFG, BPMN, DECLARE, and role process descriptions.

The first design outcome of that synthesis was the decision to organize the application around an issue-centric shared state. This choice came primarily from the mined object-centric evidence, which indicated that the issue is the main coordination object through which requests are expressed, clarified, routed, acted upon, and revisited. The BPMN artifacts then contributed the stage structure of the graph, suggesting a flow from intake to contextualization, interpretation, routing, guarded handling, execution, and closure or waiting. The declarative artifacts contributed the idea that actions should not be executed unconditionally, but should instead be subject to policy checks, approval conditions, and risk-aware gating. The textual role process descriptions were finally used to populate the role-specific agents, because they gave the clearest evidence about how repository actors interpret work and what kinds of outputs they should produce.

The population of the role-specific agents is visible directly in \texttt{prompts.py}, where each executable role is defined through an explicit prompt that constrains its focus, preferences, and limitations. The \texttt{issue\_reporter} role is defined by the prompt: ``You are the Issue Reporter Agent. You focus on intake, clarification, communication, and sustained issue coordination.'' It is instructed to prefer ``summarizing issue intent and missing context,'' ``asking clarifying questions,'' ``detecting duplicates,'' ``suggesting labels, assignees, milestones,'' and ``preparing issues for downstream work,'' while avoiding ``making deep code changes directly'' and ``closing issues without strong resolution evidence.'' This prompt translates the mined issue-reporter behavior into an executable clarification agent. In the implemented app, this means that when an issue is incomplete, vague, or insufficiently actionable, the system can route it to an agent specialized in turning ambiguity into structured downstream work rather than prematurely forcing it into implementation.

The \texttt{bot} role is specified by the prompt: ``You are the Workflow Bot Agent. You automate safe, repetitive workflow actions around issues.'' It is guided toward ``applying labels,'' ``posting status comments,'' ``requesting reviewers,'' ``reminding stalled participants,'' and ``routing issues to the next stage,'' while avoiding ``high-risk repository mutations'' and ``irreversible actions without approval.'' This prompt operationalizes the mined automation-oriented process behavior. It gives the application a role that is explicitly procedural and low-risk, suitable for repetitive coordination tasks rather than semantically complex engineering interpretation. The bot agent therefore acts as the app’s repository workflow maintainer, handling issues that resemble routine operational triggers.

The \texttt{implementation} role is defined by the prompt: ``You are the Implementation Agent. You translate issues and tasks into concrete implementation plans and proposed repository changes.'' It is instructed to prefer ``producing a task plan,'' ``linking intended changes to issue/task context,'' ``preserving traceability to commits, tasks, and issues,'' and ``handing off to QA, docs, or maintainer review as needed.'' The prompt further explains that it stands over an observed role family including contributor, feature developer, maintainer, devops engineer, quality engineer, and technical writer, while emphasizing that these roles are ``commit-centric and execution-oriented, but differ by specialization.'' This is a key design abstraction in the app. Rather than instantiating every implementation-adjacent mined role as a separate routable node, the app consolidates them into one executable development-focused agent that can interpret feature and change requests as implementation work while preserving traceability and handoff awareness.

The \texttt{maintainer} role is defined by the prompt: ``You are the Maintainer Approval Agent. You prioritize repository integrity, safe merges, and traceability.'' It is instructed to prefer ``validating merge readiness,'' ``confirming issue/task/commit linkage,'' and ``ensuring governance and release hygiene.'' Although this role was not always emphasized as a primary routed agent in the compact app narrative, it is important as a governance-oriented specialization. It captures the mined behavior associated with repository control, merge safety, and accountable release management. Its prompt defines it as an approval and integrity guardian rather than as a producer of implementation work.

The \texttt{quality} role is specified by the prompt: ``You are the Quality Agent. You focus on validation, test coverage, regression risk, and acceptance criteria.'' It is instructed to prefer ``test planning,'' ``identifying missing checks,'' ``recommending validation steps,'' and ``blocking closure when evidence is weak.'' This prompt directly encodes the mined quality-engineer behavior into an executable validation agent. In the application, this allows quality concerns to be handled as a distinct functional pathway rather than being absorbed into general development logic. The quality agent therefore represents the part of the mined process concerned with correctness, evidence, and trustworthy closure conditions.

The \texttt{technical\_writer} role is defined by the prompt: ``You are the Technical Writer Agent. You focus on documentation completeness, release-note readiness, and user-facing clarity.'' It is instructed to prefer ``docs updates,'' ``issue/PR summaries,'' ``release communication,'' and ``onboarding or usage explanations.'' This prompt gives the app a dedicated communication and documentation specialization. It is particularly important because it enables the graph to treat documentation work as first-class repository work rather than as an afterthought of implementation. When an issue is primarily about improving explanation, release communication, or onboarding clarity, this agent provides a role-consistent interpretation and corresponding action proposal.

Taken together, these prompt definitions show how the LangGraph agentic application was populated: not merely by naming roles, but by encoding mined behavioral expectations into concrete LLM instructions. Each agent in the app was specified as an executable interpretation frame with a bounded responsibility, a preferred action space, and explicit avoidance criteria. The surrounding LangGraph nodes---issue ingestion, issue context loading, intent classification, policy guarding, role routing, execution, and closure handling---provide the process skeleton, while the prompts provide the role semantics. This combination of graph structure and prompt-defined role specialization is what transformed mined repository knowledge into an operational agentic application.

\section{Evaluation}

The application was evaluated through two GitHub issues, \#1 and \#2, as a compact form of functional testing. The purpose of this evaluation was to test whether the populated role-specific agents and the surrounding LangGraph control flow behave coherently on realistic issue inputs. In both cases, the app ingested the issue, loaded its GitHub context, interpreted the text through LLM-based classification, routed the issue to a role-specific agent, and then derived a corresponding next action. The value of the evaluation lies less in statistical generalization and more in end-to-end verification that the populated app can distinguish issue types and produce role-appropriate responses.

Issue \#1 was interpreted by the application as a documentation-oriented request. Its content expressed a need that the app did not understand as requiring repository code modification, but rather as requiring clarification, explanatory improvement, or user-facing documentation work. On that basis, the routing logic selected the \texttt{technical\_writer\_agent}. This routing decision is important because it shows that the prompt-populated writer role was not ornamental: the app recognized a class of issue whose most suitable treatment was communication-centered rather than implementation-centered. The LLM-based analysis performed in the voice of the technical writer agent then derived a corresponding action focused on documentation work, such as updating or preparing explanatory material, summaries, or user-facing text. In functional terms, the app transformed the raw issue into a role-consistent documentation task. This confirms that the writer specialization was populated effectively enough to support differentiated downstream behavior.

Issue \#2 was used as the implementation-oriented counterpart. Its purpose was to test whether the app would recognize a request that implied development work and route it toward the implementation pathway. The observed behavior was instructive because the issue initially exposed a limitation of semantic routing. In an earlier wording, the issue did not signal implementation intent strongly enough, and the app routed it to the bot-oriented path instead. This outcome was useful rather than purely erroneous, because it demonstrated that routing is driven by the actual semantics extracted from the issue text. After the issue wording was revised to make the feature-oriented nature of the request more explicit, the app reinterpreted it and routed it to the \texttt{implementation\_agent}. Once there, the LLM-based analysis translated the request into a development-oriented action, treating the issue as a concrete implementation task rather than a documentation or workflow-maintenance problem. The resulting action therefore matched the population of the implementation agent as a role that converts issue statements into executable development plans and proposed repository changes.

The comparison between these two issues functions as a compact validation of the overall app design. Issue \#1 exercised the documentation pathway and showed that the technical writer prompt can drive a documentation-specific interpretation and action. Issue \#2 exercised the implementation pathway and showed that the implementation prompt can drive a development-specific interpretation and action when the issue text clearly signals feature intent. At the same time, the initial misrouting of the earlier phrasing of issue \#2 exposed a realistic characteristic of LLM-mediated issue routing: interpretation quality depends on the clarity of the issue description. This is not simply a system defect, but a reflection of the repository process itself, where ambiguous issue statements often require clarification before responsible action can be taken. In that respect, the evaluation also indirectly justifies the presence of clarification-oriented role behavior in the app.

Overall, the two-issue test demonstrates that the LangGraph application is a functioning synthesis of mined process artifacts and prompt-specified role behavior. It can consume real GitHub issues, interpret them with an LLM, route them to differentiated role agents, and derive actions that remain semantically aligned with those roles. The evaluation therefore supports the claim that the app is not merely a descriptive visualization of mined repository behavior, but an executable operationalization of that behavior in an issue-driven GitHub setting.

\end{document}

% Made with Bob
